% UD7 — Eserciziario
% PROTOTIPO: pochi esercizi scritti a mano, solo per vedere l'impaginazione.
% Il formato vero degli esercizi lo decide il ticket "Fonte unica degli esercizi".
\eserciziario{Campione del template: Lambert-Beer e radioattività}

\begin{tcolorbox}[fisbox={formula}{formulaBg},title={Come usare l'eserciziario}]
\textbf{Livello 1}: applicazione diretta di una formula.\\
\textbf{Livello 2}: due passaggi, il risultato del primo serve nel secondo.\\
\textbf{Livello 3}: problemi che combinano più concetti, spesso in un contesto medico.\\
\textbf{Quesiti d'esame}: nel formato della prova, da risolvere senza calcolatrice.

Prova da solo, poi controlla nei risultati rapidi; se il risultato non torna, leggi la soluzione svolta.
\end{tcolorbox}

\section*{Problemi}

\livello{1}{Applicazione diretta delle formule}

\problema\label{pb:ud7-hvl-tessuto}
Un tessuto ha coefficiente di attenuazione $\mu = \SI[per-mode=power]{0,35}{\centi\metre\tothe{-1}}$ per i raggi X di una radiografia. Calcola lo spessore emivalente.

\livello{2}{Problemi a due passaggi}

\problema\label{pb:ud7-piombo}
Per certi raggi $\gamma$ lo spessore emivalente del piombo è \SI{0,50}{\centi\metre}.
(a) Quale spessore riduce l'intensità a $1/16$?
(b) Qual è lo spessore minimo, in numero intero di spessori emivalenti, che riduce l'intensità sotto l'\SI{1}{\percent}?

\livello{3}{Problemi di media difficoltà}

\problema\label{pb:ud7-tecnezio}
Per una scintigrafia un paziente deve ricevere alle 14:00 un radiofarmaco al tecnezio-99m (emivita \SI{6,0}{\hour}) con attività \SI{400}{\mega\becquerel}. La dose viene preparata alle 8:00 dello stesso giorno.
(a) Quale attività deve avere alla preparazione?
(b) Quanti nuclei di tecnezio-99m riceve il paziente?
(c) Trascurando l'eliminazione biologica, quale attività resta nel paziente 24 ore dopo la somministrazione?

\section*{Quesiti d'esame}

\quesito\label{qe:ud7-beta}
Il fosforo-32, usato in alcune terapie, decade $\beta^-$. Quale nucleo si forma?
\opzioni{$\nuclide{32}{14}{Si}$}{$\nuclide{32}{16}{S}$}{$\nuclide{31}{15}{P}$}{$\nuclide{28}{13}{Al}$}{$\nuclide{32}{15}{P}$}

\quesito\label{qe:ud7-assorbanza}
Una soluzione lascia passare il \SI{10}{\percent} della luce che la colpisce. La sua assorbanza vale \completa.

\section*{Risultati rapidi}

\begin{center}
\begin{tabular}{@{}ll@{}}
\toprule
Esercizio & Risultato \\
\midrule
Problema~\ref{pb:ud7-hvl-tessuto} & $x_{1/2} \approx \SI{2,0}{\centi\metre}$ \\
Problema~\ref{pb:ud7-piombo} & (a) \SI{2,0}{\centi\metre}; (b) \SI{3,5}{\centi\metre} \\
Problema~\ref{pb:ud7-tecnezio} & (a) \SI{800}{\mega\becquerel}; (b) \num{1,3e13}; (c) \SI{25}{\mega\becquerel} \\
Quesito~\ref{qe:ud7-beta} & B \\
Quesito~\ref{qe:ud7-assorbanza} & 1 \\
Quesito lampo~\ref{ql:ud7-lambda} & A \\
Quesito lampo~\ref{ql:ud7-spettro} & C \\
Quesito lampo~\ref{ql:ud7-fotone} & 2 \\
Quesito lampo~\ref{ql:ud7-ionizzanti} & E \\
Quesito lampo~\ref{ql:ud7-hvl} & D \\
Quesito lampo~\ref{ql:ud7-alfa} & 234 \\
Quesito lampo~\ref{ql:ud7-emivita} & B \\
Quesito lampo~\ref{ql:ud7-indice} & 400 \\
Quesito lampo~\ref{ql:ud7-lente} & C \\
\bottomrule
\end{tabular}
\end{center}

\section*{Soluzioni svolte}

\subsection*{Problemi}

\solproblema{pb:ud7-hvl-tessuto}
Lo spessore emivalente è quello per cui $e^{-\mu x} = 1/2$:
\[
  x_{1/2} = \frac{\ln 2}{\mu} = \frac{0{,}69}{\SI[per-mode=power]{0,35}{\centi\metre\tothe{-1}}} \approx \SI{2,0}{\centi\metre}
\]

\solproblema{pb:ud7-piombo}
(a) $1/16 = 1/2^4$: servono 4 spessori emivalenti, cioè $4 \cdot \SI{0,50}{\centi\metre} = \SI{2,0}{\centi\metre}$.

(b) L'intensità deve scendere sotto $1/100$. Con 6 spessori emivalenti passa $1/2^6 = 1/64$, ancora più dell'\SI{1}{\percent}; con 7 passa $1/128$, meno dell'\SI{1}{\percent}. Lo spessore è $7 \cdot \SI{0,50}{\centi\metre} = \SI{3,5}{\centi\metre}$.

\solproblema{pb:ud7-tecnezio}
(a) Dalle 8:00 alle 14:00 passano \SI{6}{\hour}, cioè un'emivita: l'attività si dimezza. Alla preparazione serve il doppio, \SI{800}{\mega\becquerel}.

(b) Dalla relazione $\mathcal{A} = \lambda N$, con l'emivita in secondi ($\SI{6,0}{\hour} = \SI{21600}{\second}$):
\[
  \lambda = \frac{\ln 2}{T_{1/2}} = \frac{0{,}69}{\SI{21600}{\second}} \approx \SI[per-mode=power]{3,2e-5}{\second\tothe{-1}}
  \qquad
  N = \frac{\mathcal{A}}{\lambda} = \frac{\SI[per-mode=power]{4,0e8}{\second\tothe{-1}}}{\SI[per-mode=power]{3,2e-5}{\second\tothe{-1}}} \approx \num{1,3e13}
\]

(c) \SI{24}{\hour} sono 4 emivite: $\SI{400}{\mega\becquerel} / 2^4 = \SI{25}{\mega\becquerel}$. In realtà l'attività sarà ancora più bassa, perché una parte del radiofarmaco viene eliminata con le urine.

\subsection*{Quesiti d'esame}

\solquesito{qe:ud7-beta}
\textbf{Risposta B.} Nel decadimento $\beta^-$ un neutrone diventa protone ed esce un elettrone: il numero di massa resta 32, il numero atomico sale da 15 a 16, che è lo zolfo.
\begin{itemize}
  \item A) $Z$ diminuisce di 1: è il decadimento $\beta^+$, non $\beta^-$.
  \item C) Toglie un neutrone al nucleo: nessun decadimento $\beta$ cambia il numero di massa.
  \item D) $A - 4$ e $Z - 2$: è il decadimento $\alpha$.
  \item E) Il nucleo resta lo stesso: succede solo nel decadimento $\gamma$.
\end{itemize}

\solquesito{qe:ud7-assorbanza}
\textbf{Risposta: 1.} Se passa il \SI{10}{\percent}, $I_0/I = 10$ e $A = \log_{10} 10 = 1$.

\subsection*{Quesiti lampo}

\soluzionelampo{ql:ud7-lambda}
\textbf{Risposta A.} Nel vuoto tutte le onde elettromagnetiche hanno la stessa velocità $c = \lambda f$. Se $f$ raddoppia e $c$ resta uguale, $\lambda$ si dimezza.
\begin{itemize}
  \item B) Tratta $\lambda$ come inversamente proporzionale al quadrato della frequenza.
  \item C) Confonde la lunghezza d'onda con la velocità, che nel vuoto è la stessa per tutte le onde.
  \item D) Tratta $\lambda$ come direttamente proporzionale a $f$.
  \item E) Sbaglia due volte: proporzionalità diretta e al quadrato.
\end{itemize}

\soluzionelampo{ql:ud7-spettro}
\textbf{Risposta C.} In ordine di frequenza crescente lo spettro è radio, microonde, infrarosso, visibile, ultravioletto, X, $\gamma$; la lunghezza d'onda cresce nell'ordine opposto: $\gamma$, X, visibile, infrarosso, microonde.
\begin{itemize}
  \item A) Mette i raggi $\gamma$ fra l'ultravioletto e il visibile.
  \item B) È l'ordine di lunghezza d'onda \emph{decrescente}: scambia frequenza e lunghezza d'onda.
  \item D) Mette l'ultravioletto fra il visibile e l'infrarosso.
  \item E) Mette le microonde prima delle onde radio e salta da visibile a $\gamma$ nel verso sbagliato.
\end{itemize}

\soluzionelampo{ql:ud7-fotone}
\textbf{Risposta: 2.} Con la formula pratica $E \approx 1240/\lambda = 1240/620 = \SI{2}{\electronvolt}$.

\soluzionelampo{ql:ud7-ionizzanti}
\textbf{Risposta E.} Una radiazione è ionizzante se il singolo fotone ha un'energia maggiore dell'energia di ionizzazione, circa \SI{10}{\electronvolt}. Un fotone di microonde ha circa \SI{e-5}{\electronvolt}, troppo poco, anche se i fotoni sono moltissimi.
\begin{itemize}
  \item A) Confonde l'intensità del fascio con l'energia del singolo fotone.
  \item B) Una lunghezza d'onda lunga vuol dire fotoni di energia \emph{bassa}.
  \item C) Essere assorbita non vuol dire ionizzare: i fotoni visibili hanno circa \SIrange{2}{3}{\electronvolt}.
  \item D) I raggi X sono ionizzanti anche se non hanno carica: conta l'energia del fotone.
\end{itemize}

\soluzionelampo{ql:ud7-hvl}
\textbf{Risposta D.} Ridurre a $1/4 = 1/2^2$ vuol dire attraversare 2 spessori emivalenti, quindi $x_{1/2} = \SI{1}{\centi\metre}$. \SI{5}{\centi\metre} sono 5 spessori emivalenti: $I/I_0 = 1/2^5 = 1/32$.
\begin{itemize}
  \item A) Tratta l'attenuazione come una proporzione: «$1/4$ in \SI{2}{\centi\metre}, quindi $1/10$ in \SI{5}{\centi\metre}». L'attenuazione è esponenziale, non lineare.
  \item B) È la frazione dopo \SI{4}{\centi\metre}: manca l'ultimo centimetro.
  \item C) Applica la legge dell'inverso del quadrato, che dipende dalla distanza dalla sorgente, non dal materiale.
  \item E) Scambia $1/4$ con $1/2$ e prende \SI{2}{\centi\metre} come spessore emivalente: $1/2^{2,5} \approx 1/6$.
\end{itemize}

\soluzionelampo{ql:ud7-alfa}
\textbf{Risposta: 234.} La particella $\alpha$ porta via 4 nucleoni: $A = 238 - 4 = 234$. Il numero atomico scende da 92 a 90 (torio):
$\nuclide{238}{92}{U} \to \nuclide{234}{90}{Th} + \nuclide{4}{2}{He}$.

\soluzionelampo{ql:ud7-emivita}
\textbf{Risposta B.} Da 800 a \SI{100}{\kilo\becquerel} l'attività si riduce di un fattore $8 = 2^3$: sono passate 3 emivite. $T_{1/2} = \SI{15}{giorni}/3 = \SI{5}{giorni}$.
\begin{itemize}
  \item A) Divide il tempo per il fattore di riduzione, 8, invece che per il numero di emivite, 3.
  \item C) Pensa che l'emivita sia sempre metà del tempo trascorso.
  \item D) Considera tutto l'intervallo come una sola emivita: in un'emivita l'attività si dimezza soltanto.
  \item E) Moltiplica invece di dividere.
\end{itemize}

\soluzionelampo{ql:ud7-indice}
\textbf{Risposta: 400.} Passando nel vetro la frequenza non cambia e la velocità diventa $c/n$, quindi anche la lunghezza d'onda si divide per $n$: $\lambda = \SI{600}{\nano\metre}/1{,}5 = \SI{400}{\nano\metre}$.

\soluzionelampo{ql:ud7-lente}
\textbf{Risposta C.} L'oggetto è a $p = 2f$. Dalla legge dei punti coniugati $1/q = 1/10 - 1/20 = 1/20$, quindi $q = \SI{20}{\centi\metre}$: l'immagine è reale, oltre la lente. L'ingrandimento è $m = -q/p = -1$: capovolta e alta come l'oggetto.
\begin{itemize}
  \item A) L'immagine è virtuale solo se l'oggetto è fra la lente e il fuoco.
  \item B) Mette l'immagine nel fuoco, come se l'oggetto fosse molto lontano.
  \item D) Con una sola lente convergente un'immagine reale è sempre capovolta, e qui $|m| = 1$.
  \item E) L'immagine va all'infinito solo quando l'oggetto è nel fuoco.
\end{itemize}
